\chapter{Vitamins and minerals a z}
\index{Vitamins and minerals a z}

\section*{Definition and Overview}
Vitamins and minerals, collectively referred to as micronutrients, are essential organic and inorganic compounds required by the human body in minute quantities to facilitate a vast array of physiological processes. Unlike macronutrients (carbohydrates, proteins, and fats), micronutrients do not directly provide energy (calories) but act as vital catalysts, structural components, and regulators of metabolism. Vitamins are organic substances synthesized by plants or animals, classified based on their solubility into water-soluble (the B-complex vitamins and vitamin C) and fat-soluble (vitamins A, D, E, and K). Water-soluble vitamins are generally not stored in the body in significant amounts (with the exception of vitamin B12) and require regular dietary replenishment, as excess quantities are excreted in the urine. Conversely, fat-soluble vitamins are absorbed alongside dietary lipids and stored in the liver and adipose tissues, rendering them susceptible to accumulation and potential toxicity (hypervitaminosis) if consumed in excessive amounts.

Minerals are inorganic elements originating from the earth, soil, and water, which enter the food chain through plant uptake or animal consumption. They are categorized into macrominerals (such as calcium, phosphorus, magnesium, sodium, potassium, and chloride), which the body requires in amounts exceeding 100 mg per day, and trace minerals (including iron, zinc, iodine, selenium, copper, manganese, fluoride, and chromium), needed in microgram or milligram quantities. These micronutrients serve diverse and critical functions: calcium and phosphorus provide structural integrity to bones and teeth; sodium and potassium maintain cellular membrane potentials and fluid balance; iron is the central component of hemoglobin for oxygen transport; and vitamins function as coenzymes or cofactors in biochemical pathways, such as energy production, DNA synthesis, and cellular repair. Understanding the spectrum of vitamins and minerals is of paramount clinical and public health significance, as deficiencies and toxicities represent a major global burden of disease, often referred to as ``hidden hunger'' when subclinical.

\section*{Epidemiology}
The epidemiological landscape of micronutrient deficiencies is vast, affecting billions of individuals worldwide across all developmental and economic strata. The World Health Organization (WHO) estimates that over two billion people suffer from one or more micronutrient deficiencies, with the highest prevalence observed in low- and middle-income countries (LMICs). However, subclinical deficiencies and toxicities are also increasingly recognized in high-income nations due to dietary shifts, processed foods, and the widespread, unregulated use of dietary supplements.

Iron deficiency is the most prevalent nutritional disorder globally, affecting approximately 30\% of the world's population, which translates to over two billion individuals. It is the primary cause of anemia, particularly impacting preschool-aged children (prevalence exceeding 40\%) and pregnant women (prevalence approximately 38\%) due to their high physiological iron demands. Vitamin D deficiency is a major public health concern across all age groups, with estimates suggesting that over 50\% of the global population has insufficient vitamin D levels, and approximately 1 billion people suffer from clinical deficiency. This is primarily attributed to modern sedentary indoor lifestyles, sun-avoidance behaviors, and geographic location.

Vitamin A deficiency remains a leading cause of preventable childhood blindness, particularly in Sub-Saharan Africa and South Asia, affecting an estimated 190 million preschool-aged children. Iodine deficiency, though significantly reduced by universal salt iodization programs, still affects nearly 30\% of school-aged children globally, remaining the single most common cause of preventable brain damage and cognitive impairment. In high-income countries, specific subpopulations exhibit distinct epidemiological risks; for instance, vitamin B12 deficiency is prevalent in up to 20\% of adults over the age of 60 due to age-related atrophic gastritis, and is increasingly observed in vegetarians and vegans who do not supplement their diets.

\section*{Aetiology and Pathophysiology}
The development of micronutrient imbalances—whether deficiencies or toxicities—arises from a complex interplay of dietary habits, physiological demands, pathological processes affecting absorption, and drug-nutrient interactions.

\begin{itemize}
    \item \textbf{Inadequate Dietary Intake:} The primary cause of micronutrient deficiency is a diet lacking diversity and density. Poverty, food insecurity, and reliance on single staple crops (such as polished rice or maize) lead to widespread deficiencies in iron, zinc, and vitamin A. Restrictive dietary choices, such as strict veganism without supplementation, directly lead to cobalamin (vitamin B12) deficiency because this vitamin is synthesized exclusively by microorganisms and is found naturally only in animal-source foods.
    \item \textbf{Malabsorption Syndromes:} Gastrointestinal diseases disrupt the mechanical and chemical processes of nutrient absorption. Conditions such as celiac disease, Crohn's disease, chronic pancreatitis, and short bowel syndrome impair the mucosal surface or enzymatic breakdown necessary for absorption. For example, fat malabsorption leads to deficiencies in fat-soluble vitamins (A, D, E, and K), while atrophic gastritis and autoimmune pernicious anemia destroy gastric parietal cells, preventing the secretion of intrinsic factor required for vitamin B12 absorption.
    \item \textbf{Increased Physiological Demand:} During rapid life-stage transitions, the body's metabolic requirements for specific micronutrients rise sharply. Pregnancy and lactation increase the demands for folate (to support rapid cell division and prevent neural tube defects), iron (for expanding maternal red cell mass and fetal development), and calcium. Infants, young children, and adolescents also have elevated requirements to support bone growth and cognitive development.
    \item \textbf{Excessive Physiological Loss:} Chronic or acute loss of bodily fluids and tissues can deplete micronutrient stores. Chronic blood loss from gastrointestinal bleeding, heavy menstruation, or parasitic infections (such as hookworm) is a major driver of iron deficiency anemia. Similarly, excessive renal excretion due to chronic kidney disease or poorly controlled diabetes can lead to depletion of water-soluble vitamins and magnesium.
    \item \textbf{Drug-Nutrient Interactions:} Long-term pharmacotherapy can interfere with micronutrient absorption, metabolism, or excretion. Proton pump inhibitors (PPIs) raise gastric pH, impairing the cleavage of vitamin B12 from dietary protein and reducing calcium and magnesium absorption. Anticonvulsants like phenytoin accelerate the hepatic metabolism of vitamin D, and methotrexate acts as a direct folate antagonist by inhibiting dihydrofolate reductase.
    \item \textbf{Hypervitaminosis and Mineral Toxicity:} Nutrient excess is almost exclusively caused by the overconsumption of high-dose dietary supplements rather than dietary intake. Because fat-soluble vitamins are stored in the body, excessive intake of vitamin A can saturate binding proteins, leading to free retinoid toxicity that damages cell membranes and the liver. High doses of vitamin D lead to excessive intestinal calcium absorption and bone resorption, culminating in hypercalcemia.
\end{itemize}

\section*{Clinical Features}
The clinical manifestations of micronutrient disorders range from subtle, non-specific constitutional symptoms to severe, life-threatening clinical syndromes. The signs and symptoms depend heavily on the specific nutrient involved, the severity of the depletion, and the chronicity of the state.

\begin{itemize}
    \item \textbf{Vitamin A Deficiency:} Early features include night blindness (nyctalopia) due to impaired regeneration of rhodopsin in the retina. Progression leads to xerophthalmia, characterized by conjunctival dryness, Bitot's spots (foamy keratin plaques on the bulbar conjunctiva), corneal ulceration, and ultimately keratomalacia (liquefactive necrosis of the cornea). Skin manifestations include follicular hyperkeratosis (phrynoderma).
    \item \textbf{Thiamine (Vitamin B1) Deficiency:} Presents as wet beriberi, dry beriberi, or Wernicke-Korsakoff syndrome. Wet beriberi is characterized by high-output cardiac failure, peripheral edema, tachycardia, and cardiomegaly. Dry beriberi presents as a symmetrical peripheral neuropathy with sensory loss, muscle wasting, and loss of deep tendon reflexes. Wernicke's encephalopathy is an acute medical emergency presenting with the triad of ophthalmoplegia/nystagmus, ataxia, and confusion, which can progress to Korsakoff's psychosis (marked by anterograde amnesia and confabulation).
    \item \textbf{Niacin (Vitamin B3) Deficiency:} Causes pellagra, clinically defined by the ``four Ds'': dermatitis, diarrhea, dementia, and death. The dermatitis is photosensitive, presenting with symmetrical, hyperpigmented, scaling rashes on sun-exposed areas, classically forming a ring around the neck (Casal's necklace).
    \item \textbf{Vitamin B12 and Folate Deficiency:} Both cause megaloblastic anemia, presenting with fatigue, pallor, dyspnea, and glossitis (a smooth, red, painful tongue). However, vitamin B12 deficiency uniquely causes subacute combined degeneration of the spinal cord due to impaired myelin synthesis in the posterior and lateral columns, presenting with symmetric paresthesias, loss of vibratory and proprioceptive sensation, ataxia, and progressive spastic paraparesis.
    \item \textbf{Vitamin C Deficiency (Scurvy):} Manifests due to impaired collagen synthesis, presenting with follicular hyperkeratosis, coiled ``corkscrew'' hairs, perifollicular hemorrhages, easy bruising (petechiae and ecchymoses), swollen and bleeding gums, loose teeth, joint pain (hemarthrosis), and delayed wound healing.
    \item \textbf{Vitamin D Deficiency:} In children, it causes rickets, characterized by skeletal deformities (bowed legs, knock-knees, rachitic rosary of the ribs, craniotabes). In adults, it presents as osteomalacia, characterized by diffuse bone pain, muscle weakness, and a waddling gait.
    \item \textbf{Iron Deficiency:} Clinically manifests as fatigue, pale conjunctivae and skin, dyspnea, cold extremities, pica (an abnormal craving for non-nutritive substances like ice, clay, or starch), pagophagia, and koilonychia (spoon-shaped nails).
    \item \textbf{Calcium Deficiency:} Acute hypocalcemia causes neuromuscular irritability, presenting with circumoral paresthesias, muscle cramps, tetany, and positive Chvostek's sign (facial muscle twitching upon tapping the facial nerve) and Trousseau's sign (carpopedal spasm induced by inflating a blood pressure cuff).
    \item \textbf{Zinc Deficiency:} Characterized by growth retardation, delayed sexual maturation, alopecia, diarrhea, impaired wound healing, immune dysfunction leading to recurrent infections, dysgeusia (loss of taste), and a distinctive skin rash known as acrodermatitis enteropathica.
\end{itemize}

\section*{Differential Diagnosis}
Because the symptoms of micronutrient disorders are often non-specific (such as fatigue, neuropathy, and skin changes), it is essential to differentiate them from other systemic, hematological, and neurological diseases.

\begin{itemize}
    \item \textbf{Myelodysplastic Syndromes (MDS):} MDS frequently presents with macrocytic anemia, leukopenia, and thrombocytopenia, closely mimicking the hematological profile of vitamin B12 or folate deficiency. A bone marrow biopsy is crucial to differentiate the clonal dysplastic changes of MDS from the megaloblastic changes of vitamin deficiency.
    \item \textbf{Lead Poisoning:} Lead toxicity impairs heme synthesis, resulting in a microcytic, hypochromic anemia with basophilic stippling, similar to iron deficiency anemia. Blood lead level measurements and evaluation for lead-induced neuropathy help distinguish these conditions.
    \item \textbf{Diabetic Neuropathy:} The distal, symmetric, sensory neuropathy of diabetes can mimic dry beriberi or the early stages of subacute combined degeneration of the spinal cord. Detailed glycemic history, HbA1c testing, and assessment of spinal cord tract involvement (proprioception and vibration, which are spared in pure diabetic neuropathy but impaired in B12 deficiency) are necessary.
    \item \textbf{Hypothyroidism:} Presents with fatigue, cold intolerance, dry skin, hair loss, and cognitive slowing, mimicking iodine, selenium, or thyroid hormone-related disorders. Serum thyroid-stimulating hormone (TSH) and free thyroxine (T4) levels are diagnostic.
    \item \textbf{Autoimmune Connective Tissue Diseases:} Conditions like systemic lupus erythematosus (SLE) or vasculitis can present with skin lesions, mucosal bleeding, joint pain, and constitutional symptoms, resembling scurvy. Serum autoantibody profiles (ANA, anti-dsDNA) and inflammatory markers (ESR, CRP) help differentiate these systemic inflammatory processes.
    \item \textbf{Primary Hyperparathyroidism:} Causes hypercalcemia, bone resorption, and renal stones, which can be misdiagnosed as vitamin D toxicity. Measuring parathyroid hormone (PTH) levels is key; PTH is elevated in primary hyperparathyroidism but suppressed in vitamin D toxicity.
\end{itemize}

\section*{When to Seek Medical Care}
While minor micronutrient deficiencies can be managed through gradual dietary adjustments and standard oral supplements, certain acute presentations represent severe physiological imbalances that demand urgent medical evaluation or emergency intervention.

Patients must seek immediate emergency medical care (by calling local emergency services in many countries, 911 in the United States, or 999 in the United Kingdom) if they experience any of the following life-threatening symptoms:
\begin{itemize}
    \item Acute, severe shortness of breath, chest pain, or a rapid, irregular heartbeat, which may indicate severe anemia or high-output heart failure (wet beriberi).
    \item Sudden, progressive neurological changes, including confusion, disorientation, severe loss of physical coordination, unsteady gait, or paralysis, which are signs of Wernicke's encephalopathy or advanced subacute combined degeneration.
    \item Severe, uncontrollable muscle spasms, seizures, or throat tightness (laryngospasm), which are signs of acute, severe hypocalcemia (tetany).
    \item Accidental or intentional ingestion of massive amounts of iron supplements, particularly in young children, which can cause acute, life-threatening iron poisoning.
\end{itemize}

A primary care physician or specialist should be consulted promptly for persistent, non-emergency symptoms such as chronic fatigue, progressive numbness or tingling in the hands or feet, unexplained bone pain, slow-healing wounds, severe hair loss, recurrent infections, or vision changes in low light.

\section*{Diagnosis and Investigations}
A structured diagnostic approach combining a detailed dietary history, physical examination, and targeted laboratory and imaging investigations is necessary to confirm micronutrient imbalances and identify their underlying causes.

\begin{itemize}
    \item \textbf{Clinical Examination:} A thorough physical examination should evaluate the skin (for rashes, petechiae, follicular hyperkeratosis), eyes (for dryness, Bitot's spots, conjunctival pallor), mouth (for glossitis, angular cheilitis, bleeding gums), and neurological system (testing reflexes, vibration sense, proprioception, and gait).
    \item \textbf{Hematological Investigations:} A Complete Blood Count (CBC) and peripheral blood smear are fundamental. Iron deficiency classically demonstrates a microcytic, hypochromic anemia (low Mean Corpuscular Volume [MCV] and low Mean Corpuscular Hemoglobin Concentration [MCHC]). Vitamin B12 and folate deficiencies demonstrate a macrocytic anemia (high MCV) with characteristic hypersegmented neutrophils on the peripheral smear.
    \item \textbf{Biochemical Assays:} Direct serum measurements are widely used. Iron status is assessed using a complete iron panel, including serum ferritin (the most sensitive indicator of iron stores), serum iron, Total Iron-Binding Capacity (TIBC), and transferrin saturation. Vitamin D status is determined by measuring serum 25-hydroxyvitamin D [25(OH)D], which is the major circulating form. Serum cobalamin (B12) and red blood cell folate levels assess B12 and folate status.
    \item \textbf{Functional Biomarkers:} When serum levels are borderline, functional biomarkers are measured. Methylmalonic acid (MMA) and homocysteine are elevated in B12 deficiency, whereas only homocysteine is elevated in folate deficiency. Thiamine status can be evaluated functionally by measuring red blood cell transketolase activity.
    \item \textbf{Bone and Mineral Biomarkers:} For calcium and vitamin D disorders, investigations include serum total and ionized calcium, inorganic phosphate, alkaline phosphatase (ALP), and parathyroid hormone (PTH).
    \item \textbf{Imaging and Densitometry:} Skeletal radiographs are indicated in children with suspected rickets to identify cupping and fraying of the metaphyses. In adults, Dual-Energy X-ray Absorptiometry (DEXA) scans are utilized to measure bone mineral density (BMD) at the femoral neck and lumbar spine, diagnosing osteopenia or osteoporosis resulting from chronic calcium and vitamin D depletion.
\end{itemize}

\section*{Management}
The primary objective of management is to replenish depleted micronutrient stores safely, alleviate clinical symptoms, and address the underlying cause of the imbalance to prevent recurrence.

\begin{itemize}
    \item \textbf{Dietary Interventions:} The foundation of managing subclinical and mild deficiencies is dietary modification. Patients are advised to consume a diverse diet rich in micronutrients. This includes leafy green vegetables (rich in folate and calcium), lean red meats and organ meats (rich in bioavailable heme iron and vitamin B12), seafood and iodized salt (for iodine), nuts, seeds, and whole grains (for zinc and magnesium). Fortified foods, such as milk fortified with vitamin D or cereals fortified with iron and B vitamins, are valuable resources.
    \item \textbf{Oral Pharmacological Supplementation:} Clinical deficiencies require therapeutic doses of oral supplements. Iron deficiency anemia is typically treated with oral ferrous sulfate, ferrous fumarate, or ferrous gluconate, often co-administered with vitamin C to enhance intestinal absorption. Vitamin D deficiency is corrected with high-dose cholecalciferol (vitamin D3) or ergocalciferol (vitamin D2). Oral folate is administered daily to correct folate deficiency.
    \item \textbf{Parenteral and Intravenous Therapy:} Indicated when oral absorption is severely compromised or rapid replenishment is clinically critical. Intramuscular injections of cyanocobalamin or hydroxocobalamin are mandatory for patients with pernicious anemia or significant ileal resection. Intravenous iron infusions (e.g., iron carboxymaltose or ferric derisomaltose) are preferred for patients with severe iron deficiency, intolerance to oral iron, or active inflammatory bowel disease. Acute Wernicke's encephalopathy is treated with immediate high-dose intravenous thiamine before any glucose administration to prevent irreversible brain damage.
    \item \textbf{Addressing the Primary Pathology:} Long-term resolution requires treating the underlying cause. This involves implementing a strict gluten-free diet in patients with celiac disease, administering anti-inflammatory therapies to control Crohn's disease, or prescribing pancreatic enzyme replacement therapy for pancreatic insufficiency.
    \item \textbf{Treatment of Toxicity:} Management of hypervitaminosis involves immediate discontinuation of the offending supplement, supportive care, and measures to lower serum levels. For vitamin D toxicity with severe hypercalcemia, treatment includes aggressive intravenous hydration with normal saline, loop diuretics to promote urinary calcium excretion, and bisphosphonates or calcitonin to inhibit bone resorption.
\end{itemize}

\section*{Complications}
Neglecting the diagnosis and treatment of micronutrient deficiencies or toxicities can result in permanent structural damage, severe organ dysfunction, and increased mortality.

\begin{itemize}
    \item \textbf{Irreversible Neurological Damage:} Untreated vitamin B12 deficiency leads to progressive, permanent demyelination of the spinal cord (subacute combined degeneration), culminating in permanent ataxia, loss of motor function, and spastic paraplegia. Prolonged thiamine deficiency in Wernicke's encephalopathy can progress to Korsakoff's psychosis, characterized by permanent, irreversible memory impairment.
    \item \textbf{Cardiovascular Compromise:} Chronic severe anemia (from iron, B12, or folate deficiency) and wet beriberi place a high hemodynamic load on the heart, leading to high-output cardiac failure, cardiomegaly, and fatal arrhythmias.
    \item \textbf{Blindness and Visual Impairment:} Severe vitamin A deficiency can lead to keratomalacia, causing rapid corneal melting, perforation, scarring, and irreversible bilateral blindness.
    \item \textbf{Developmental and Congenital Anomalies:} Maternal folate deficiency during the first weeks of gestation increases the risk of neural tube defects, such as spina bifida and anencephaly. Maternal iodine deficiency leads to congenital hypothyroidism (formerly termed cretinism), resulting in permanent mental retardation, deaf-mutism, and short stature in the offspring.
    \item \textbf{Severe Skeletal Deformities and Fractures:} Chronic calcium and vitamin D deficiency lead to progressive bone loss, resulting in permanent skeletal deformities in children and an elevated risk of fragility fractures (especially hip and vertebral fractures) in adults.
    \item \textbf{Systemic Toxicity Complications:} Chronic iron overload (hemochromatosis or hemosiderosis from repeated transfusions or excessive supplementation) causes iron deposition in parenchymal organs, leading to hepatic cirrhosis, cardiomyopathy, diabetes mellitus, and skin hyperpigmentation (bronze diabetes). Severe hypercalcemia from vitamin D toxicity can lead to nephrocalcinosis, renal calculi, and acute or chronic kidney failure.
\end{itemize}

\section*{Prognosis and Outcomes}
The prognosis for individuals with micronutrient imbalances is generally excellent, provided the condition is diagnosed early and managed appropriately. Most nutritional deficiencies respond rapidly and completely to targeted supplementation. For example, in patients with megaloblastic or microcytic anemia, hematological parameters begin to normalize within one to two weeks, with hemoglobin levels returning to baseline within two to three months. Clinical symptoms such as fatigue, pica, and glossitis typically resolve within days of initiating therapy.

However, the prognosis is heavily influenced by the duration and severity of the deficiency prior to intervention. If structural or neurological damage has occurred, the recovery may be incomplete. Neurological deficits from vitamin B12 and thiamine deficiencies may improve slowly over several months, but long-standing, severe deficits often leave residual impairment. Similarly, bone deformities caused by pediatric rickets may persist into adulthood, requiring orthopedic correction. In cases of acute toxicity, such as severe iron poisoning or vitamin D-induced hypercalcemia, the prognosis depends on the speed of emergency medical intervention to prevent organ damage. Compliance with long-term dietary recommendations and maintenance supplementation, particularly in patients with irreversible malabsorptive conditions or strict dietary restrictions, is the primary determinant of long-term success.

\section*{Prevention and Self-Care}
Preventing micronutrient imbalances relies on maintaining a balanced, nutrient-dense diet, incorporating targeted supplementation for vulnerable populations, and raising public health awareness regarding safe dietary practices.

\begin{itemize}
    \item \textbf{Balanced and Diverse Diet:} The most effective way to prevent deficiencies is to eat a wide variety of foods. The diet should include at least five portions of fruits and vegetables daily, whole grains, lean proteins (poultry, fish, eggs, legumes), and healthy fats. Incorporating fermented foods and dairy products helps maintain adequate calcium and B-vitamin intake.
    \item \textbf{Targeted Supplementation:} Certain life stages and dietary patterns require proactive supplementation. All women planning a pregnancy or who are of childbearing age should take 400 micrograms of folic acid daily to prevent neural tube defects. Strict vegans and vegetarians must consume B12-fortified foods or take a regular vitamin B12 supplement. Breastfed infants in dark-skinned populations or northern latitudes may require vitamin D supplementation.
    \item \textbf{Safe Sun Exposure:} Exposure of the arms and legs to direct sunlight for 10--15 minutes daily (avoiding peak UV times to minimize skin cancer risk) helps maintain adequate endogenous vitamin D synthesis.
    \item \textbf{Avoidance of High-Dose Supplementation:} Individuals should avoid self-prescribing high-dose single-nutrient supplements without medical supervision. Adhering to the Recommended Dietary Allowance (RDA) and not exceeding the Tolerable Upper Intake Level (UL) prevents toxicities.
    \item \textbf{Public Health Fortification Programs:} Utilizing fortified consumer goods, such as iodized salt, vitamin D-fortified milk, and folic acid-fortified wheat flour, provides a baseline level of protection against common deficiencies.
    \item \textbf{Routine Monitoring:} Individuals who have undergone bariatric surgery, have chronic gastrointestinal diseases, or take long-term medications (like PPIs or metformin) should undergo annual blood tests to monitor their micronutrient status and adjust their treatment plans accordingly.
\end{itemize}

\section*{Sources and Further Reading}
\begin{itemize}
    \item World Health Organization (WHO) - Micronutrients: \url{https://www.who.int/health-topics/micronutrients}
    \item National Institutes of Health (NIH) Office of Dietary Supplements: \url{https://ods.od.nih.gov}
    \item Centers for Disease Control and Prevention (CDC) - Micronutrient Facts: \url{https://www.cdc.gov/nutrition/micronutrient-facts/index.html}
    \item Harvard T.H. Chan School of Public Health - The Nutrition Source: \url{https://www.hsph.harvard.edu/nutritionsource}
    \item British Dietetic Association (BDA) - Food Fact Sheets: \url{https://www.bda.uk.com/resource-library.html}
    \item Dietitians Australia - Nutrition Information: \url{https://dietitiansaustralia.org.au/health-advice}
    \item Mayo Clinic - Drugs and Supplements Reference: \url{https://www.mayoclinic.org/drugs-supplements}
\end{itemize}